\chapter{الذرات المتعددة الإلكترونات ورموز الحدود}\label{ch:b3:many-electron-atoms}

في سنة 1868 كشف خط أصفر في طيف الشمس عن عنصر
مجهول على الأرض، هو الهيليوم. وحين عُزل أخيرًا في المختبر، سنة
1895، بدا طيفه كأنه طيف \emph{عنصرين}: عائلتان من
الخطوط، نُسبتا حينًا إلى «أورثوهيليوم»
و«باراهيليوم»، ولم تتبادلا الضوء قط. وليس
ثمة إلا هيليوم واحد. فالعائلتان هما الحالات الثلاثية والحالات الأحادية للذرة
نفسها، يفصل بينهما سبين الإلكترون وقاعدة أعمق
من أي قوة: يجب أن تغيّر الدالة الموجية لإلكترونين إشارتها حين
يتبادلان موضعيهما. يتتبع هذا الفصل تلك القاعدة من مبدأ باولي إلى
رموز الحدود التي يسمّي بها أهل المطيافية كل حالة لكل ذرة
وكل أيون --- بما فيها أيونات الفلزات الانتقالية التي تُفسَّر ألوانها
في \cref{ch:b3:complex-spectra-magnetism}.

\begin{recall}
وسم مجلد السنة الأولى الإلكترونات بأربعة أعداد كمية $n$ و$l$
و$m_l$ و$m_s$، وبنى التوزيعات الأساسية بمبدأ باولي 
وقاعدة كليتشكوفسكي وقاعدة هوند، وعدّ الإلكترونات المنفردة. وأعطى مجلد السنة الثانية
طاقات المدارات وقواعد سلاتر للحجب.
وقدّم \cref{ch:b3:quantum-model-systems} المؤثرات والقيم الذاتية وكمية
الحركة الزاوية للدوّار ($\hat L^2$ بقيم ذاتية $l(l+1)\hbar^2$) 
وذرة الهيدروجين.
\end{recall}

\section{السبين والمدارات المغزلية}

يحمل الإلكترون كمية حركة زاوية ذاتية، هي السبين، عدده
الكمي $s = \frac12$: فللمؤثر $\hat S^2$ \omterm{def:b3:quantum-model-systems:operator}{قيمة ذاتية} وحيدة $s(s+1)\hbar^2 =
\frac34\hbar^2$، وللمؤثر $\hat S_z$ القيمتان الذاتيتان $m_s\hbar = \pm\frac12\hbar$.
وتُكتب دالتا السبين $\alpha$ ($m_s = +\frac12$) و$\beta$
($m_s = -\frac12$)؛ وهما متعامدتان ومنظَّمتان، $\langle\alpha|\alpha\rangle =
\langle\beta|\beta\rangle = 1$ و$\langle\alpha|\beta\rangle = 0$. وليس للسبين
نظير كلاسيكي؛ فهو ليس دوران كرة صغيرة، ولا
يدخل البتة في \omterm{def:b3:quantum-model-systems:schrodinger}{مؤثر هاملتون} غير النسبوي.

\begin{definition}[المدار المغزلي]\label{def:b3:many-electron-atoms:spin-orbital}
\emph{المدار المغزلي}\index{مدار مغزلي} جداء مدار فضائي
$\phi(\mathbf r)$ ودالة سبين: $\phi\alpha$ أو $\phi\beta$. 
والمدار المغزلي يحمل الوصف الكامل لإلكترون واحد.
\end{definition}

فالمدار الذري الموسوم $(n, l, m_l)$ في مجلد السنة الأولى يعطي إذن \omterm{def:b3:many-electron-atoms:spin-orbital}{مدارين
مغزليين}، ولهذا تتسع الطبقة الفرعية لعدد $2(2l+1)$ من الإلكترونات. والسؤال
الحقيقي هو لماذا لا يتسع \omterm{def:b3:many-electron-atoms:spin-orbital}{المدار المغزلي} إلا لإلكترون واحد على الأكثر.

\section{الإلكترونات غير القابلة للتمييز ومحدد سلاتر}

\begin{definition}[الجسيمات غير القابلة للتمييز، الدالة الموجية ضد المتناظرة]\label{def:b3:many-electron-atoms:indistinguishable}
تكون الجسيمات \emph{غير قابلة للتمييز}\index{جسيمات غير قابلة للتمييز} إذا
لم يستطع أي قياس أن يميّز بينها: فكل إلكترون مطابق لكل إلكترون
آخر. وتكون الدالة الموجية لعدة إلكترونات \emph{دالة موجية ضد
متناظرة}\index{دالة موجية ضد متناظرة} إذا غيّرت إشارتها حين تتبادل
إحداثيات (الفضاء والسبين) أي إلكترونين:
$\Psi(\dots,i,\dots,j,\dots) = -\Psi(\dots,j,\dots,i,\dots)$.
\end{definition}

عدم القابلية للتمييز وحده يقتضي ألا تتغير $|\Psi|^2$ بالتبادل،
فلا يسع $\Psi$ إلا أن تغيّر إشارتها على الأكثر. وقد اختارت الطبيعة:
فالدالة الموجية لأي نظام من الإلكترونات --- لأي فرميونات --- ضد
متناظرة. وهذه مسلّمة من مسلّمات ميكانيكا الكم، يؤكدها كل طيف
ذري.

\begin{definition}[محدد سلاتر]\label{def:b3:many-electron-atoms:slater}
من أجل $N$ إلكترونًا في $N$ \omterm{def:b3:many-electron-atoms:spin-orbital}{مدارًا مغزليًا} مختلفًا $\chi_1, \dots, \chi_N$، يكون
\emph{محدد سلاتر}\index{محدد سلاتر}
\[
\Psi = \frac{1}{\sqrt{N!}}\begin{vmatrix}
\chi_1(1) & \chi_2(1) & \cdots & \chi_N(1)\\
\chi_1(2) & \chi_2(2) & \cdots & \chi_N(2)\\
\vdots & & & \vdots\\
\chi_1(N) & \chi_2(N) & \cdots & \chi_N(N)
\end{vmatrix},
\]
حيث يحمل السطر $i$ الإلكترون $i$ في كل \omterm{def:b3:many-electron-atoms:spin-orbital}{مدار مغزلي}. ويُكتب اختصارًا
$|\chi_1\chi_2\dots\chi_N|$.
\end{definition}

\begin{theorem}[مبدأ باولي]\label{thm:b3:many-electron-atoms:pauli}
\omterm{def:b3:many-electron-atoms:slater}{محدد سلاتر} ضد متناظر، وينعدم إذا شغل إلكترونان
\omterm{def:b3:many-electron-atoms:spin-orbital}{المدار المغزلي} نفسه. ومنه لا يمكن أن يكون لإلكترونين في ذرة واحدة
الأعداد الكمية الأربعة نفسها.
\end{theorem}

\begin{proof}
تبادل الإلكترونين $i$ و$j$ يبادل سطرين من المحدد،
فتتغير إشارته: فضد التناظر محقق لكل زوج. وإذا تساوى
\omterm{def:b3:many-electron-atoms:spin-orbital}{مداران مغزليان}، تساوى عمودان وانعدم المحدد:
فلا وجود لمثل هذه الحالة. والإلكترونان في الذرة نفسها اللذان لهما $n$ و$l$
و$m_l$ و$m_s$ نفسها يقعان في \omterm{def:b3:many-electron-atoms:spin-orbital}{المدار المغزلي} نفسه.
\end{proof}

\begin{example}[الحالة الأساسية للهيليوم]\label{ex:b3:many-electron-atoms:helium-ground}
والإلكترونان كلاهما في $1s$،
\[
\Psi = \frac{1}{\sqrt2}\begin{vmatrix}1s\alpha(1) & 1s\beta(1)\\ 1s\alpha(2) &
1s\beta(2)\end{vmatrix} = 1s(1)\,1s(2)\cdot\frac{1}{\sqrt2}\big[\alpha(1)\beta(2)
- \beta(1)\alpha(2)\big].
\]
الجزء الفضائي متناظر، وجزء السبين ضد متناظر: فالسبينان
مزدوجان، والسبين الكلي معدوم.
\end{example}

\section{الهيليوم: التنافر الكولومي والتبادل}

لنُثِر أحد إلكتروني الهيليوم إلى $2s$. يمكن بناء أربعة محددات من
$1s\alpha$ و$1s\beta$ و$2s\alpha$ و$2s\beta$ بإلكترون واحد في كل مدار.
وبإعادة ترتيبها، تصير جداءات دالة فضائية ودالة سبين:
\[
\Psi_\pm = \frac{1}{\sqrt2}\big[1s(1)2s(2) \pm 2s(1)1s(2)\big]\times(\text{دالة
سبين}).
\]
يجب أن تُقرن الدالة الفضائية ضد المتناظرة $\Psi_-$ بدالة سبين
متناظرة، وهي ثلاث: $\alpha(1)\alpha(2)$ 
و$\beta(1)\beta(2)$ و$\frac{1}{\sqrt2}[\alpha(1)\beta(2) + \beta(1)\alpha(2)]$، أي
المركبات الثلاث $M_S = 1, 0, -1$ لسبين كلي $S = 1$. ويجب أن تُقرن الدالة المتناظرة
$\Psi_+$ بدالة السبين ضد المتناظرة الوحيدة، $S = 0$.

\begin{definition}[الحالتان الأحادية والثلاثية]\label{def:b3:many-electron-atoms:singlet-triplet}
الحالة ذات السبين الكلي $S = 0$ \emph{حالة أحادية}\index{حالة أحادية}؛
والحالة ذات السبين الكلي $S = 1$، التي لمركباتها الثلاث $M_S = 1, 0, -1$
الطاقة نفسها في غياب مجال مغناطيسي، \emph{حالة
ثلاثية}\index{حالة ثلاثية}.
\end{definition}

\begin{definition}[تكامل التبادل]\label{def:b3:many-electron-atoms:exchange}
من أجل مدارين $a$ و$b$ والتنافر الإلكتروني $e^2/4\pi\varepsilon_0
r_{12}$، يكون \emph{تكامل التبادل}\index{تكامل التبادل}
\[
K_{ab} = \iint a(1)b(2)\,\frac{e^2}{4\pi\varepsilon_0r_{12}}\,b(1)a(2)\,\dd\tau_1\dd\tau_2,
\]
وهو التكامل نفسه الذي يعطي التنافر الإلكتروني الكلاسيكي $J_{ab}$ (مع
$a(1)b(2)$ في الطرفين) غير أن الإلكترونين متبادلان في
أحد الطرفين. و$K_{ab}$ موجب وليس له نظير كلاسيكي.
\end{definition}

\begin{proposition}[طاقتا الحالتين الأحادية والثلاثية]\label{prop:b3:many-electron-atoms:helium}
في الرتبة الأولى بالنسبة إلى التنافر الإلكتروني، تكون طاقات الحالات $1s2s$ للهيليوم
$E_\pm = E_{1s} + E_{2s} + J \pm K$، \omterm{def:b3:many-electron-atoms:singlet-triplet}{والحالة الأحادية} ($+$) أعلى من
\omterm{def:b3:many-electron-atoms:singlet-triplet}{الحالة الثلاثية} ($-$) بمقدار $2K$.
\end{proposition}

\begin{proof}
الطاقة غير المضطربة للدالتين $\Psi_\pm$ كلتيهما هي $E_{1s} + E_{2s}$ (مداران شبيهان
بمدارات الهيدروجين بشحنة نووية 2). وتصحيح الرتبة الأولى هو
\omterm{def:b3:quantum-model-systems:hermitian}{القيمة المتوقعة} للتنافر $\hat V_{12}$ في الحالة المنظَّمة:
$\frac12\langle 1s(1)2s(2) \pm 2s(1)1s(2)|\hat V_{12}|1s(1)2s(2) \pm
2s(1)1s(2)\rangle$. يعطي الحدان المباشران $\frac12(J + J)$؛ ويعطي الحدان
المتقاطعان $\pm\frac12(K + K)$، لأن $\hat V_{12}$ متناظر بالنسبة إلى
الإلكترونين. ومنه $\Delta E = J \pm K$، أما دوال السبين، المنظَّمة
التي لا يمسّها $\hat V_{12}$، فلا تضرب إلا في 1.
\end{proof}

تقع \omterm{def:b3:many-electron-atoms:singlet-triplet}{الحالة الثلاثية} أدنى مع أن \omterm{def:b3:quantum-model-systems:schrodinger}{مؤثر هاملتون} لا يحتوي على السبين: فدالتها
الفضائية تنعدم حين $\mathbf r_1 = \mathbf r_2$، فيتجنب الإلكترونان
أحدهما الآخر ويتنافران أقل. وهذا «ثقب فيرمي» هو
المضمون الفيزيائي لقاعدة هوند الأولى.

\begin{example}[تكامل التبادل للهيليوم، مقيسًا]\label{ex:b3:many-electron-atoms:K}
% ledger: lev:He.2s3S1, lev:He.2s1S0
تقع مستويات $1s2s$ للهيليوم عند \qty{159856}{cm^{-1}} (\babelsublr{${}^3$S}) 
و\qty{166277}{cm^{-1}} (\babelsublr{${}^1$S}) فوق الحالة الأساسية: \omterm{def:b3:many-electron-atoms:singlet-triplet}{فالحالة الأحادية} أعلى بمقدار
\qty{6421}{cm^{-1}}، أي \qty{0.796}{eV}، و$K(1s,2s) = \qty{0.398}{eV}$.
ويبيّن الشكل أدناه العائلتين.
\end{example}

\section{اقتران راسل--ساوندرز ورموز الحدود}

في الذرة الخفيفة، تتجمع كميات الحركة الزاوية المدارية للإلكترونات في
كمية كلية $\mathbf L$، وسبيناتها في كمية كلية $\mathbf S$، ثم لا يقترن
$\mathbf L$ و$\mathbf S$ أحدهما بالآخر إلا اقترانًا ضعيفًا.

\begin{definition}[اقتران راسل--ساوندرز]\label{def:b3:many-electron-atoms:russell-saunders}
في \emph{اقتران راسل--ساوندرز}\index{اقتران راسل--ساوندرز}، تُوسم
حالات التوزيع بواسطة \emph{العدد الكمي لكمية الحركة الزاوية المدارية
الكلية}\index{العدد الكمي لكمية الحركة الزاوية المدارية الكلية}
$L$ ($\hat{\mathbf L}^2 = L(L+1)\hbar^2$، مع $M_L = \sum m_l$)، و\emph{العدد الكمي
للسبين الكلي}\index{العدد الكمي للسبين الكلي} $S$ ($M_S = \sum m_s$)،
و\emph{العدد الكمي لكمية الحركة الزاوية الكلية}\index{العدد الكمي لكمية الحركة الزاوية الكلية}
$J$، الذي يأخذ القيم $|L - S|, \dots, L + S$.
\end{definition}

\begin{definition}[الحد الطيفي، تعددية السبين، رمز الحد]\label{def:b3:many-electron-atoms:term}
\emph{الحد الطيفي}\index{حد طيفي} هو مجموعة الحالات
$(2L+1)(2S+1)$ لتوزيع ما ذات $L$ و$S$ معطيين. 
و\emph{تعددية السبين}\index{تعددية السبين} له هي $2S + 1$. و\emph{رمز الحد}
\index{رمز الحد} هو ${}^{2S+1}L$، مع الحروف S وP وD وF وG وH
للقيم $L = 0, 1, 2, 3, 4, 5$؛ وتُكتب الحالة ذات $J$ المعطى ${}^{2S+1}L_J$ ---
مثلًا \termsym{3}{P}{2}.
\end{definition}

\begin{proposition}[الطبقات المكتملة]\label{prop:b3:many-electron-atoms:closed-shell}
للطبقة الفرعية الممتلئة $L = 0$ و$S = 0$؛ فهي لا تساهم بشيء في
\omterm{def:b3:many-electron-atoms:term}{رمز الحد}.
\end{proposition}

\begin{proof}
في الطبقة الفرعية الممتلئة تظهر كل قيمة $m_l$ مرتين وتظهر كل قيمة $m_s = \pm\frac12$
$2l+1$ مرة، ومنه $M_L = 0$ و$M_S = 0$؛ وليس ثمة إلا طريقة واحدة
لملئها (محدد واحد). والحالة الوحيدة ذات $M_L = M_S = 0$ لا يمكن أن
تنتمي إلا إلى $L = 0$، $S = 0$.
\end{proof}

\begin{proposition}[عدّ حالات التوزيع]\label{prop:b3:many-electron-atoms:count}
تعطي $N$ إلكترونات في طبقة فرعية ذات $2(2l+1)$ \omterm{def:b3:many-electron-atoms:spin-orbital}{مدارًا مغزليًا}
$\binom{2(2l+1)}{N}$ محددًا، ومجموع \omterm{def:b3:quantum-model-systems:degenerate}{درجات الانحلال} $(2L+1)(2S+1)$
لحدودها يساوي هذا العدد.
\end{proposition}

\begin{proof}
يتعيّن المحدد بمجموعة \omterm{def:b3:many-electron-atoms:spin-orbital}{المدارات المغزلية} المشغولة (فترتيبها
لا يغيّر إلا إشارته): نختار $N$ مدارًا من بين $2(2l+1)$. والحدود هي
الحالات نفسها مُعاد تجميعها، فيجب أن يتفق العددان.
\end{proof}

\begin{method}[حدود توزيع ما]\label{met:b3:many-electron-atoms:terms}
\begin{enumerate}
  \item اكتب المحددات التي يسمح بها مبدأ باولي ورتّبها
        في جدول بحسب $M_L$ و$M_S$.
  \item جد أكبر قيمة للعدد $M_L$؛ فهي مع أكبر قيمة للعدد $M_S$ ترافقها
        تفتتح حدًا فيه $L = M_L$، $S = M_S$.
  \item اشطب محددًا واحدًا لكل زوج $(M_L, M_S)$ يحقق
        $|M_L| \le L$، $|M_S| \le S$.
  \item كرّر مع ما تبقى، حتى يفرغ الجدول؛ وتحقق من العدد.
\end{enumerate}
\end{method}

\begin{example}[حدود $p^2$]\label{ex:b3:many-electron-atoms:p2}
يعطي إلكترونان في $p$ ($m_l = 1, 0, -1$) عددًا من المحددات قدره $\binom62 = 15$.
وجدولها بحسب $M_L$ و$M_S$ هو:
\begin{center}\small
\begin{tabular}{c|ccc}
\hline
 & $M_S = 1$ & $M_S = 0$ & $M_S = -1$\\ \hline
$M_L = 2$ & -- & $(1^+,1^-)$ & --\\
$M_L = 1$ & $(1^+,0^+)$ & $(1^+,0^-)$, $(1^-,0^+)$ & $(1^-,0^-)$\\
$M_L = 0$ & $(1^+,-1^+)$ & $(1^+,-1^-)$, $(1^-,-1^+)$, $(0^+,0^-)$ & $(1^-,-1^-)$\\
$M_L = -1$ & $(0^+,-1^+)$ & $(0^+,-1^-)$, $(0^-,-1^+)$ & $(0^-,-1^-)$\\
$M_L = -2$ & -- & $(-1^+,-1^-)$ & --\\ \hline
\end{tabular}
\end{center}
(حيث $m_l^\pm$ من أجل $m_s = \pm\frac12$). لا تظهر $M_L = 2$ إلا مع $M_S = 0$:
فهي حد \babelsublr{${}^1$D} (5 حالات). وأكبر قيمة متبقية $M_L = 1$ تأتي مع
$M_S = 1$: حد \babelsublr{${}^3$P} (9 حالات). وتبقى حالة واحدة ذات $M_L = M_S = 0$: حد
\babelsublr{${}^1$S}. وبالفعل $5 + 9 + 1 = 15$.
\end{example}

\begin{example}[الكربون، مقيسًا]\label{ex:b3:many-electron-atoms:carbon}
% ledger: lev:C.3P1, lev:C.3P2, lev:C.1D2, lev:C.1S0
يعطي التوزيع الأساسي $2p^2$ للكربون المستويات \termsym{3}{P}{0}
و\termsym{3}{P}{1} و\termsym{3}{P}{2} عند 0 و16.4 و\qty{43.4}{cm^{-1}}، ثم
\termsym{1}{D}{2} عند \qty{10193}{cm^{-1}} و\termsym{1}{S}{0} عند
\qty{21648}{cm^{-1}}: الحد \babelsublr{${}^3$P} هو الأدنى، كما تتنبأ قواعد هوند.
\end{example}

\begin{omfigure}
\begin{tikzpicture}
\begin{axis}[name=main, width=0.62\linewidth, height=6.4cm, xmin=-0.3, xmax=6.6,
  ymin=-1500, ymax=23500, axis y line=left, axis x line=none,
  ylabel={\foreignlanguage{arabic}{الطاقة \babelsublr{(cm$^{-1}$)}}}, unbounded coords=jump, clip=false,
  scaled y ticks=false, ytick={0,10000,20000},
  tick label style={font=\small}, label style={font=\small}]
\addplot[omstate] table[x=x, y=E] {figdata/out/bachelor-3-atomic-levels-carbon.dat};
\draw[densely dotted, gray] (axis cs:1,4858.5) -- (axis cs:2,29.6);
\draw[densely dotted, gray] (axis cs:1,4858.5) -- (axis cs:2,10192.7);
\draw[densely dotted, gray] (axis cs:1,4858.5) -- (axis cs:2,21648);
\draw[densely dotted, gray] (axis cs:3,10192.7) -- (axis cs:4,10192.7);
\draw[densely dotted, gray] (axis cs:3,21648) -- (axis cs:4,21648);
\node[font=\small, anchor=south] at (axis cs:0.5,4858.5) {$2p^2$};
\node[font=\small, anchor=south] at (axis cs:2.5,21648) {${}^1$S};
\node[font=\small, anchor=south] at (axis cs:2.5,10192.7) {${}^1$D};
\node[font=\small, anchor=south] at (axis cs:2.5,29.6) {${}^3$P};
\node[font=\small, anchor=west] at (axis cs:5.05,21648) {\termsym{1}{S}{0}};
\node[font=\small, anchor=west] at (axis cs:5.05,10192.7) {\termsym{1}{D}{2}};
\node[font=\scriptsize, anchor=north] at (axis cs:0.5,-600) {\foreignlanguage{arabic}{التوزيع}};
\node[font=\scriptsize, anchor=north] at (axis cs:2.5,-600) {\foreignlanguage{arabic}{الحدود}};
\node[font=\scriptsize, anchor=north] at (axis cs:4.5,-600) {\foreignlanguage{arabic}{المستويات}};
\draw[gray, ->] (axis cs:3.05,60) -- (axis cs:6.9,2500);
\end{axis}
\begin{axis}[at={(main.south east)}, anchor=south west, xshift=1.2cm, yshift=0.9cm,
  width=3.6cm, height=3.6cm, xmin=-0.1, xmax=1.1, ymin=-5, ymax=50,
  axis y line=right, axis x line=none, unbounded coords=jump, ytick={0,16.4,43.4},
  yticklabels={0,16.4,43.4}, tick label style={font=\scriptsize}, clip=false,
  title={\scriptsize \foreignlanguage{arabic}{\babelsublr{${}^3$P} مكبَّرًا}}, title style={yshift=-4pt}]
\addplot[omstate] table[x=x, y=E] {figdata/out/bachelor-3-atomic-levels-carbon-zoom.dat};
\node[font=\scriptsize, anchor=east] at (axis cs:0,0) {$J = 0$};
\node[font=\scriptsize, anchor=east] at (axis cs:0,16.4) {$J = 1$};
\node[font=\scriptsize, anchor=east] at (axis cs:0,43.4) {$J = 2$};
\end{axis}
\end{tikzpicture}

{\small التوزيع الأساسي للكربون: توزيع واحد، وثلاثة حدود
(عند طاقاتها المقيسة؛ الحد \babelsublr{${}^3$P} عند المتوسط المرجَّح
لمستوياته، والتوزيع عند متوسط الحالات الخمس عشرة كلها) وخمسة
مستويات. وانشطار السبين--المدار للحد \babelsublr{${}^3$P} (المكبَّر، بوحدة
\unit{cm^{-1}}) أصغر ببضع مئات المرات من الفواصل بين
الحدود.}
\end{omfigure}

\section{قواعد هوند واقتران السبين--المدار}

\begin{proposition}[قواعد هوند للحدود]\label{prop:b3:many-electron-atoms:hund-terms}
من أجل التوزيع الأساسي لذرة أو أيون، يكون الحد الأدنى هو الحد
ذا (1) أكبر $S$، ثم (2) من أجل هذه القيمة للعدد $S$، أكبر $L$. وأدنى
مستوياته له (3) $J = |L - S|$ إذا كانت الطبقة الفرعية أقل من نصف ممتلئة، و$J = L +
S$ إذا كانت أكثر من نصف ممتلئة.
\end{proposition}
\begin{proof}[مكانتها]
هذه القواعد مثبتة تجريبيًا، ومن أجل التوزيعات الأساسية وحدها؛
وليست مبرهنات. فالقاعدة 1 هي استقرار التبادل الذي رأيناه في القسم
السابق؛ والقاعدة 2 تقول إن الإلكترونات التي تدور في الاتجاه نفسه تلتقي
أقل؛ والقاعدة 3 تنتج عن إشارة \omterm{def:b3:many-electron-atoms:spin-orbit}{اقتران السبين--المدار}، أدناه.
\end{proof}

\begin{method}[الحد الأساسي من مخطط الخانات]\label{met:b3:many-electron-atoms:ground-term}
\begin{enumerate}
  \item املأ خانات الطبقة الفرعية المفتوحة بدءًا من $m_l = +l$ نزولًا، بإلكترون
        واحد لكل خانة بسبينات متوازية، ثم زاوج بدءًا من $+l$ من جديد.
  \item $S = \frac12 \times$(عدد الإلكترونات المنفردة)؛ $L = |\sum m_l|$.
  \item $J = |L - S|$ إذا كانت أقل من نصف ممتلئة، و$L + S$ إذا كانت أكثر من نصف
        ممتلئة؛ وإذا كانت نصف ممتلئة تمامًا فإن $L = 0$ و$J = S$.
\end{enumerate}
\end{method}

\begin{example}[الحدود الأساسية]\label{ex:b3:many-electron-atoms:ground-terms}
% ledger: lev:Ti2.3P0, lev:Ni2.3P0, lev:Cr3.4P1_2
\ce{Ti^2+}، $d^2$: الخانتان $m_l = 2, 1$ مملوءتان، $S = 1$، $L = 3$، $J = 2$:
\termsym{3}{F}{2}. \ce{Cr^3+}، $d^3$: $S = \frac32$، $L = 2 + 1 + 0 = 3$،
\termsym{4}{F}{3/2}. \ce{Fe^2+}، $d^6$: خمسة إلكترونات للأعلى، وواحد للأسفل في $m_l = 2$:
$S = 2$، $L = 2$، أكثر من نصف ممتلئة، \termsym{5}{D}{4}. \ce{Ni^2+}، $d^8$:
$S = 1$، $L = 3$، \termsym{3}{F}{4}. وكلها توافق المستوى الأساسي الذي
تقيسه المطيافية الذرية.
\end{example}

\begin{definition}[اقتران السبين--المدار، البنية الدقيقة]\label{def:b3:many-electron-atoms:spin-orbit}
\emph{اقتران السبين--المدار}\index{اقتران السبين--المدار} هو التآثر
بين العزمين المغناطيسيين لسبين الإلكترون ولحركته
المدارية، ويُكتب $A\,\hat{\mathbf L}\cdot\hat{\mathbf S}/\hbar^2$ من أجل حد. وهو
يشطر الحد إلى مستوياته ذات قيم $J$ المختلفة: وهذه \emph{البنية
الدقيقة}\index{البنية الدقيقة} للحد.
\end{definition}

\begin{proposition}[قاعدة فواصل لانديه]\label{prop:b3:many-electron-atoms:lande}
داخل الحد، $E(J) = E_0 + \frac{A}{2}[J(J+1) - L(L+1) - S(S+1)]$، ومنه
$E(J) - E(J - 1) = AJ$. و$A > 0$ (عادية) للطبقة الفرعية الأقل من نصف ممتلئة،
و$A < 0$ (مقلوبة) للأكثر من نصف ممتلئة.
\end{proposition}

\begin{proof}
تعطي $\hat{\mathbf J} = \hat{\mathbf L} + \hat{\mathbf S}$ العلاقة $\hat{\mathbf J}^2 =
\hat{\mathbf L}^2 + \hat{\mathbf S}^2 + 2\hat{\mathbf L}\cdot\hat{\mathbf S}$، ومنه
$\hat{\mathbf L}\cdot\hat{\mathbf S} = \frac12(\hat{\mathbf J}^2 - \hat{\mathbf L}^2 -
\hat{\mathbf S}^2)$، وقيمته الذاتية في حالة ذات $J$ و$L$ و$S$ معطاة هي
$\frac{\hbar^2}{2}[J(J+1) - L(L+1) - S(S+1)]$. والفرق بين $J$ 
و$J - 1$ هو $\frac A2[J(J+1) - (J-1)J] = AJ$. أما إشارة $A$ فنقبلها دون برهان: فالطبقة
الفرعية الأكثر من نصف ممتلئة تسلك سلوك العدد الموافق من الثقوب
الموجبة.
\end{proof}

\begin{example}[اختبار قاعدة الفواصل]\label{ex:b3:many-electron-atoms:lande-test}
% ledger: lev:C.3P1, lev:C.3P2, lev:O.3P1, lev:O.3P0
من أجل \babelsublr{${}^3$P} تتنبأ القاعدة بفاصلين نسبتهما $J = 2 : 1$، أي 2.
الكربون: $(43.4 - 16.4)/16.4 = 1.65$. الأكسجين ($2p^4$، مقلوب، \termsym{3}{P}{2}
الأدنى، ثم $J = 1$ عند \qty{158.3}{cm^{-1}} و$J = 0$ عند \qty{227.0}{cm^{-1}}):
$158.3/68.7 = 2.30$. فالقاعدة صحيحة في حدود 20\,\%: \omterm{def:b3:many-electron-atoms:russell-saunders}{اقتران راسل--ساوندرز}
وصف جيد للذرات الخفيفة، لا وصف دقيق.
\end{example}

\begin{omfigure}
\begin{tikzpicture}[font=\small]
  \draw[omstate] (0,0) -- (1.6,0) node[midway, below] {$3s\ {}^2$S$_{1/2}$};
  \draw[omstate] (0,3.0) -- (1.6,3.0);
  \node[anchor=east] at (-0.1,3.0) {$3p$};
  \draw[omstate] (2.8,2.75) -- (4.4,2.75) node[right] {${}^2$P$_{1/2}$};
  \draw[omstate] (2.8,3.25) -- (4.4,3.25) node[right] {${}^2$P$_{3/2}$};
  \draw[densely dotted, gray] (1.6,3.0) -- (2.8,2.75) (1.6,3.0) -- (2.8,3.25);
  \draw[omfluo] (3.3,2.73) -- (3.3,0.02) node[pos=0.55, left] {D$_1$};
  \draw[omfluo] (3.9,3.23) -- (3.9,0.02) node[pos=0.55, right] {D$_2$};
  \draw[omstate] (2.8,0) -- (4.4,0);
  \draw[densely dotted, gray] (1.6,0) -- (2.8,0);
  \draw[decorate, decoration={brace, amplitude=3pt}] (5.5,3.27) -- (5.5,2.73)
     node[midway, right=3pt, align=left] {\footnotesize \qty{17.2}{cm^{-1}}\\[-1pt]\footnotesize \foreignlanguage{arabic}{دون مراعاة المقياس}};
\end{tikzpicture}

{\small خطا الصوديوم D. ينشطر المستوى $3p$ \omterm{def:b3:many-electron-atoms:spin-orbit}{باقتران السبين--المدار}
إلى \babelsublr{${}^2$P$_{1/2}$} و\babelsublr{${}^2$P$_{3/2}$}؛ ويعطي الإصدار نحو المستوى الأساسي $3s$
خطين أصفرين، \babelsublr{D$_1$} عند \qty{589.76}{nm} و\babelsublr{D$_2$} عند \qty{589.16}{nm} (طولا
الموجتين في الفراغ).}
\end{omfigure}

\begin{example}[ثنائية الصوديوم]\label{ex:b3:many-electron-atoms:sodium}
% ledger: lev:Na.3p1_2, lev:Na.3p3_2
يقع المستويان $3p$ للصوديوم عند 16956.17 و\qty{16973.37}{cm^{-1}}: فخطا D
يقعان عند $10^7/16956.17 = \qty{589.76}{nm}$ و\qty{589.16}{nm} في الفراغ،
ويفصل بينهما \qty{17.20}{cm^{-1}}. ومن أجل \babelsublr{${}^2$P}، $E(\frac32) - E(\frac12) =
\frac32A$، ومنه $A = \qty{11.47}{cm^{-1}}$.
\end{example}

\begin{proposition}[قواعد الانتقاء للذرات]\label{prop:b3:many-electron-atoms:selection-atoms}
في \omterm{def:b3:many-electron-atoms:russell-saunders}{اقتران راسل--ساوندرز}، يقتضي الانتقال ثنائي القطب الكهربائي
$\Delta S = 0$، $\Delta L = 0, \pm1$، $\Delta J = 0, \pm1$ (لكن ليس $J = 0 \to J
= 0$)، وتغيّر الزوجية: من أجل قفزة إلكترون واحد، $\Delta l = \pm1$.
\end{proposition}
\begin{proof}[مكانتها]
تُبرهن $\Delta S = 0$ في \cref{ch:b3:electronic-spectroscopy}: \omterm{def:b3:quantum-model-systems:operator}{فمؤثر} ثنائي
القطب لا يؤثر في السبين. أما القواعد الأخرى فتنتج عن الطابع المتجهي
\omterm{def:b3:quantum-model-systems:operator}{لمؤثر} ثنائي القطب ونقبلها هنا دون برهان.
\end{proof}

\begin{omfigure}
\begin{tikzpicture}
\begin{axis}[width=0.62\linewidth, height=6.2cm, xmin=-0.5, xmax=3.5,
  ymin=158, ymax=173, axis y line=left, axis x line=none,
  ylabel={\foreignlanguage{arabic}{الطاقة \babelsublr{($10^3$ cm$^{-1}$)}}}, unbounded coords=jump, clip=false,
  tick label style={font=\small}, label style={font=\small}]
\addplot[omstate] table[x=x, y=E] {figdata/out/bachelor-3-atomic-levels-helium.dat};
\node[font=\small, anchor=west] at (axis cs:1.05,166.277) {$1s2s\ {}^1$S};
\node[font=\small, anchor=west] at (axis cs:1.05,171.135) {$1s2p\ {}^1$P};
\node[font=\small, anchor=west] at (axis cs:3.05,159.856) {$1s2s\ {}^3$S};
\node[font=\small, anchor=west] at (axis cs:3.05,169.087) {$1s2p\ {}^3$P};
\node[font=\small, anchor=south] at (axis cs:0.5,172.3) {\foreignlanguage{arabic}{الأحادية}};
\node[font=\small, anchor=south] at (axis cs:2.5,172.3) {\foreignlanguage{arabic}{الثلاثية}};
\draw[omfluo] (axis cs:0.5,171.0) -- (axis cs:0.5,166.42) node[pos=0.5, left, font=\scriptsize] {\qty{2059}{nm}};
\draw[omfluo] (axis cs:2.5,168.95) -- (axis cs:2.5,160.0) node[pos=0.5, left, font=\scriptsize] {\qty{1083}{nm}};
\draw[densely dotted, gray] (axis cs:-0.3,159.856) -- (axis cs:2,159.856);
\draw[<->, omThm] (axis cs:-0.2,159.856) -- (axis cs:-0.2,166.277) node[pos=0.5, right, font=\scriptsize] {$2K$};
\end{axis}
\end{tikzpicture}

{\small «الهيليومان». المستويات الأحادية والثلاثية للتوزيعين
$1s2s$ و$1s2p$ عند طاقاتها المقيسة فوق الحالة الأساسية (الحد
\babelsublr{${}^3$P} عند متوسط مستوياته). تصل الخطوط بين حالات العائلة
نفسها فقط؛ وكل \omterm{def:b3:many-electron-atoms:singlet-triplet}{حالة ثلاثية} تقع أدنى من نظيرتها الأحادية، بمقدار $2K$.}
\end{omfigure}

\begin{history}[الهيليومان ومبدأ الاستبعاد]
سُمّي الهيليوم باسم الشمس سنة 1868 ووجده ويليام
رامزي على الأرض سنة 1895. وحيّرت سلسلتا خطوطه أهل المطيافية ثلاثين
سنة. وفي سنة 1925 اقترح فولفغانغ باولي ألا يشترك إلكترونان في
الأعداد الكمية نفسها، وهي السنة نفسها التي اقترح فيها جورج
أولنبك وصموئيل غاودسميت سبين الإلكترون؛ وفي سنة 1926 بيّن فيرنر هايزنبرغ أن
تناظر الدالة الموجية وحده يشطر الهيليوم إلى عائلتيه الأحادية
والثلاثية.
\end{history}

\begin{inthelab}[فصل ثنائية الصوديوم]
يوضع مصباح صوديوم أمام شق مطياف ذي محزوز. ومع
محزوز ذي 600 خط في الميليمتر، تختلف زاويتا الخطين في الرتبة الأولى
بنحو \num{0.02}\textdegree: فالمطياف الجيد يفصلهما، أما
موشور الطلاب فلا. ويعمل المصباح ساخنًا؛ فلا يُلمس إلا بعد أن
يبرد.
\end{inthelab}

\section{تمارين}

\begin{exercise}[$\star$]\label{exo:b3:many-electron-atoms:1}
اكتب \omterm{def:b3:many-electron-atoms:slater}{محدد سلاتر} للحالة الأساسية لليثيوم، $1s^22s$، مع
إلكترون $2s$ ذي السبين $\alpha$. لماذا لا يوجد محدد من أجل
$1s^3$؟
\end{exercise}

\begin{exercise}[$\star$]\label{exo:b3:many-electron-atoms:2}
عُدّ محددات التوزيعات $p^2$ و$p^3$ و$d^2$. حدود
$p^3$ هي \babelsublr{${}^4$S} و\babelsublr{${}^2$D} و\babelsublr{${}^2$P}، وحدود $d^2$ هي \babelsublr{${}^3$F} و\babelsublr{${}^3$P} و\babelsublr{${}^1$G}
و\babelsublr{${}^1$D} و\babelsublr{${}^1$S}: تحقق من العددين كليهما.
\end{exercise}

\begin{exercise}[$\star$]\label{exo:b3:many-electron-atoms:3}
أعطِ الحد الأساسي والمستوى الأساسي لكل من N وO وF و\ce{Ti^2+} و\ce{Fe^3+} ($d^5$).
\end{exercise}

\begin{exercise}[$\star$]\label{exo:b3:many-electron-atoms:4}
اكتب مستويات الحدين \babelsublr{${}^2$D} و\babelsublr{${}^3$P} مع \omterm{def:b3:quantum-model-systems:degenerate}{درجات انحلالها} $2J+1$،
وتحقق من أن مجموع \omterm{def:b3:quantum-model-systems:degenerate}{درجات الانحلال} يساوي $(2L+1)(2S+1)$.
\end{exercise}

\begin{exercise}[$\star\star$]\label{exo:b3:many-electron-atoms:5}
تقع مستويات \babelsublr{${}^3$P} للكربون عند 0 و16.4 و\qty{43.4}{cm^{-1}}. احسب
نسبة الفاصلين، وثابت السبين--المدار $A$ من كل فاصل.
ماذا يدل عليه الفرق بين قيمتي $A$؟
\end{exercise}

\begin{exercise}[$\star\star$]\label{exo:b3:many-electron-atoms:6}
يقع المستويان $3p$ للصوديوم عند 16956.17 و\qty{16973.37}{cm^{-1}}. احسب
طولي موجتي خطي D في الفراغ، والفاصل بينهما بوحدة \unit{cm^{-1}} 
وبالميلي إلكترون فولت، والثابت $A$ للحد $3p$.
\end{exercise}

\begin{exercise}[$\star\star$]\label{exo:b3:many-electron-atoms:7}
أي هذه الانتقالات للصوديوم مسموح في الإصدار: $3p \to 3s$،
$3d \to 3s$، $3d \to 3p$، $4s \to 3s$، $4s \to 3p$؟ علّل كل إجابة.
\end{exercise}

\begin{exercise}[$\star\star$]\label{exo:b3:many-electron-atoms:8}
% ledger: lev:He.2p3P2, lev:He.2p3P1, lev:He.2p3P0, lev:He.2p1P1
مستويات $1s2p$ للهيليوم هي \termsym{3}{P}{2} و\termsym{3}{P}{1}
و\termsym{3}{P}{0} عند 169086.76 و169086.84 و\qty{169087.83}{cm^{-1}}، 
و\termsym{1}{P}{1} عند \qty{171134.89}{cm^{-1}}. احسب الطاقة المتوسطة
للحد \babelsublr{${}^3$P}، ثم $K(1s,2p)$ بالإلكترون فولت. قارنها بالقيمة $K(1s,2s) = \qty{0.398}{eV}$
وفسّر الفرق.
\end{exercise}

\begin{exercise}[$\star\star$]\label{exo:b3:many-electron-atoms:9}
جد حدود $d^2$ بطريقة هذا الفصل، بدءًا من
أكبر قيمة للعدد $M_L$ (يمكنك ألا تكتب الجدول كله، لكن علّل كل حد).
\end{exercise}

\begin{exercise}[$\star\star\star$]\label{exo:b3:many-electron-atoms:10}
بيّن أن الدالة الفضائية الثلاثية $\frac{1}{\sqrt2}[1s(1)2s(2) -
2s(1)1s(2)]$ تنعدم حين $\mathbf r_1 = \mathbf r_2$ وأن الدالة الأحادية
لا تنعدم. واربط ذلك بإشارة $E_+ - E_-$.
\end{exercise}

\begin{exercise}[$\star\star\star$]\label{exo:b3:many-electron-atoms:11}
% ledger: lev:Ni2.3F3, lev:Ni2.3F2
من أجل \ce{Ni^2+} ($d^8$) تقع مستويات \babelsublr{${}^3$F} عند 0 ($J = 4$) و1360.7 ($J = 3$)
و\qty{2269.6}{cm^{-1}} ($J = 2$). فسّر هذا الترتيب، واحسب نسبة
الفاصلين وقارنها بتنبؤ لانديه. وأعطِ $A$ من الفاصل
الأكبر.
\end{exercise}

\begin{exercise}[$\star\star\star$]\label{exo:b3:many-electron-atoms:12}
بيّن أن \omterm{def:b3:quantum-model-systems:hermitian}{القيمة المتوقعة} للمقدار $\hat{\mathbf L}\cdot\hat{\mathbf S}$، مجموعةً
على حالات الحد كلها وعددها $(2L+1)(2S+1)$، معدومة، فلا يغيّر \omterm{def:b3:many-electron-atoms:spin-orbit}{اقتران
السبين--المدار} المتوسط المرجَّح للمستويات. وتحقق من ذلك على
الحد \babelsublr{${}^3$P} للكربون.
\end{exercise}

\section{مسألة: الهيليومان}

\begin{problem}[{مسألة نهاية الأسبوع --- العائلتان الأحادية والثلاثية
للهيليوم: المحددات، وحدود أيونات الفلزات الانتقالية، وقواعد الانتقاء
التي تفصل بين العائلتين، وتكامل التبادل مقيسًا}]
\label{pb:b3:many-electron-atoms:1}
% ledger: lev:He.2s3S1, lev:He.2s1S0, lev:He.2p3P2, lev:He.2p3P1, lev:He.2p3P0, lev:He.2p1P1, lev:Ti2.3F3, lev:Ti2.3F4, lev:Ti2.3P0, lev:Ti2.3P1, lev:Ti2.3P2
المعطيات (NIST، بوحدة \unit{cm^{-1}} فوق الحالة الأساسية لكل نوع). الهيليوم:
$1s2s$ \termsym{3}{S}{1} 159855.97، \termsym{1}{S}{0} 166277.44؛ $1s2p$
\termsym{3}{P}{2,1,0} 169086.76، 169086.84، 169087.83، \termsym{1}{P}{1}
171134.89. \ce{Ti^2+}: \termsym{3}{F}{2,3,4} 0، 184.9، 420.4؛ \termsym{3}{P}{0,1,2}
10538.4، 10603.6، 10721.2. $\qty{1}{eV} = \qty{8065.54}{cm^{-1}}$.

\medskip\noindent\textbf{الجزء الأول --- \omterm{def:b3:many-electron-atoms:spin-orbital}{المدارات المغزلية} والمحددات.}
\begin{enumerate}
  \item اكتب \omterm{def:b3:many-electron-atoms:spin-orbital}{المدارات المغزلية} الأربعة المتاحة للتوزيع $1s2s$.
  \item اكتب المحددات الأربعة التي فيها إلكترون واحد في $1s$ وواحد في
        $2s$.
  \item بيّن أن $|1s\alpha\,2s\alpha|$ جداء دالة فضائية ضد متناظرة
        ودالة السبين $\alpha(1)\alpha(2)$.
  \item اجمع $|1s\alpha\,2s\beta|$ و$|1s\beta\,2s\alpha|$ للحصول على
        المركبة $M_S = 0$ للدالة الفضائية نفسها، وعلى حالة رابعة.
  \item أي دالة فضائية، المتناظرة أم ضد المتناظرة، تقترن
        بالحالة $S = 0$؟ وأيها بالحالة $S = 1$؟
  \item لماذا تسمى هاتان الحالتان أحادية وثلاثية؟
\end{enumerate}

\medskip\noindent\textbf{الجزء الثاني --- حدود \ce{Ti^2+}.}
\begin{enumerate}[resume]
  \item أعطِ توزيع \ce{Ti^2+} وعدد
        محدداته.
  \item جد حده الأساسي ومستواه الأساسي بقواعد هوند؛ وتحقق منهما
        بالمعطيات.
  \item حدوده الأخرى هي \babelsublr{${}^3$P} و\babelsublr{${}^1$G} و\babelsublr{${}^1$D} و\babelsublr{${}^1$S}: تحقق من عدد
        الحالات.
  \item هل \omterm{def:b3:many-electron-atoms:spin-orbit}{البنية الدقيقة} للحد \babelsublr{${}^3$F} عادية أم مقلوبة؟ اختبر قاعدة
        فواصل لانديه.
  \item احسب الطاقتين المتوسطتين المرجَّحتين للحدين \babelsublr{${}^3$F} و\babelsublr{${}^3$P}.
  \item الفاصل بينهما $15B$، حيث $B$ أحد وسيطي راكاه
        للأيون (\cref{ch:b3:complex-spectra-magnetism}). احسب $B$.
\end{enumerate}

\medskip\noindent\textbf{الجزء الثالث --- لماذا عائلتان.}
\begin{enumerate}[resume]
  \item أي قاعدة انتقاء تمنع الانتقال بين \omterm{def:b3:many-electron-atoms:singlet-triplet}{حالة أحادية} \omterm{def:b3:many-electron-atoms:singlet-triplet}{وحالة
        ثلاثية}؟
  \item لا تستطيع أدنى \omterm{def:b3:many-electron-atoms:singlet-triplet}{حالة ثلاثية}، $1s2s$ \termsym{3}{S}{1}، أن تُصدر نحو الحالة
        الأساسية $1s^2$ \termsym{1}{S}{0}. أعطِ سببين. وبمَ تسمى هذه
        الحالة؟
  \item احسب طول موجة الخط \babelsublr{$1s2p\ {}^3$P $\to 1s2s\ {}^3$S} (استعمل
        متوسط مستويات \babelsublr{${}^3$P}).
  \item احسب طول موجة \babelsublr{$1s2p\ {}^1$P $\to 1s2s\ {}^1$S}.
  \item احسب طول موجة \babelsublr{$1s2p\ {}^1$P $\to 1s^2\ {}^1$S}. في أي
        منطقة من الطيف يقع؟
  \item فسّر لماذا اعتقد أهل المطيافية في القرن التاسع عشر، وهم يرون هذه الخطوط،
        بوجود غازين مختلفين.
\end{enumerate}

\medskip\noindent\textbf{الجزء الرابع --- \omterm{def:b3:many-electron-atoms:exchange}{تكامل التبادل}، مقيسًا.}
\begin{enumerate}[resume]
  \item اكتب طاقتي الحالتين الأحادية والثلاثية $1s2s$ في الرتبة الأولى
        بدلالة $E_{1s}$ و$E_{2s}$ و$J$ و$K$.
  \item أيهما أدنى، ولماذا، من حيث مواضع الإلكترونين؟
  \item احسب الفجوة بين الحالتين الأحادية والثلاثية $1s2s$ بوحدة \unit{cm^{-1}} وبالإلكترون فولت.
  \item احسب $K(1s,2p)$ بالإلكترون فولت من مستويات $1s2p$.
  \item لماذا يكون $K(1s,2s)$ أكبر من $K(1s,2p)$؟
  \item هل تؤيد الحالات المثارة للهيليوم قاعدة هوند الأولى؟
  \item اذكر النتيجة: \omterm{def:b3:many-electron-atoms:exchange}{تكامل التبادل} $K(1s,2s)$ للهيليوم، بالإلكترون
        فولت.
\end{enumerate}
\end{problem}
